\section{Setting and notation}
\btodo{Work in progress}
\btodo{Decide whether we want to write ``example'', ``user'', or ``record'' for the general unit of privacy. We used ``example'' in the multi-epoch paper, so goiing with that for now.}
We build on the multi-epoch setting of \citet{choquette22multiepoch}, which we summarize here.  

Assume a database of $\numusers$ examples (users in FL, records in general DP applications) that is processed as a stream over $\dimension$ steps. A batch of $\batchsize$ examples is selected on each step $i$, and processed via an adaptively chosen function (e.g., computing a gradient at the current model), producing a vector of $\ell_2$ norm at most $\clipnorm$. These vectors are summed and provided to the DP mechanism as $\obs_i \in \R^\mdim$, which releases a privatized function of $[\obs_1, \dots, \obs_i]$, the stream so far. When a particular example contributes to the sum $\obs_i$, we say it \emph{participates} on step $i$. We are primarily interested in the case where $\numusers < \dimension \cdot \batchsize$, and hence each example is used on more than one step. This is the multiple epoch setting of ML.

\btodo{We probably need to also re-state the definition of adjacent streams. It might clarify things to introduce notation for individual examples, as well as a ``data processing function'' which assigns examples to batches for each iteration.}

\paragraph{Participation schemas}
A \emph{participation schema} $\Pi$ is a set of possible \emph{participation patterns} $\pi \in \Pi$, with each $\pi \subseteq [\dimension]$ indicating a set of steps in which any single example might participate. That is, if example $j$ participates in rounds $r_j \subseteq [n]$, then we require $r_j \subseteq \pi$ for some $\pi \in \Pi$.

Simple examples include \emph{single-participation} (each record contributes at most once, $\Pi = \big\{ \{1\}, \{2\}, \dots \{\dimension\}\big\}$) and \emph{every-step participation} (each record might contribue on every step, as in full gradient descent, $\Pi = \{ [\dimension]\}$).  \citet{choquette22multiepoch} studied \emph{fixed-epoch-order participation}, denoted \emph{\multiepoch-participation}, where each example participates at most $\maxpart$ times, with any adjacent participations exactly $\minsep$ steps apart: formally, $\Pi$ is the set of all $\pi$ such that $|\pi| \le \maxpart$, and if $\pi = \{i_1, \dots, i_k\}$ indexed in increasing order, we have $\forall j \in \{2, \dots, k\}, i_j - i_{j-1} = \minsep$. For example $\maxpartsep{3}{2}$-participation has $\Pi = \{\{1, 3, 5\}, \{2, 4, 6\}\}$. Both single and every-step participation are special cases of \multiepoch-participation, and as discussed in \citet{choquette22multiepoch}, this setting faithfully captures standard multi-epoch ML training setups.

\paragraph{Banded matrices}
We say a (general) matrix $\bfX$ is $\bands$-banded if for all $i,j \in [n]$, $|i -j| \ge \bands$ implies $\bfX\idx{i}{j} = 0$. While this is off-by-one from the typical definition of bandwidth ($\bfX$ has bandwidth $b-1$), our definition will be useful as it will be natural to match $\bands$-banded matrices with $\minsep$-min-separation. Further, for $\bands$-banded lower-triangular matrices (which will play a central role), $\bands$ intuitively gives the number of bands in the matrix.  

%\btodo{A notation issue: In our code (e.g., \texttt{sensitivity.py}, in some places we use $k$ to denote the number of bands below the main diagonal (following the use of $k$ in numpy and e.g. \url{https://en.wikipedia.org/wiki/Band_matrix}). On the other hand, we claim $k$ is max participation even in the code (and here). Proposal: we keep $k$ for max participation (updating the code). The trickier question (in the code and the paper) is what to do for min-separation and number of bands in the matrix. Strawman Proposal: Use $b$ for both, defining so that in the common case this provides the necessary orthogonality (you want a $b$-banded matrix for $b$-minsep).  This matches how we defined separation in the multi-epoch paper so that every-round participation corresponds to $b=1$. In order to ensure all participations by the same example are orthogonal, this requires a diagonal matrix, which suggests we should define a $(b=1)$-banded matrix to be a diagonal matrix (which, intuitively, also makes sense).  The only place this is awkward is with non-lower-triangular symmetric banded matrices, where it's not clear what $b=2$ would mean. (This is presumably why numpy uses $k$ the way it does). The primary place this will come up for us will be with $X = C^T C$, which when $C$ is lower-triangular banded will be a symmetric banded matrix. But (if this issue actually comes up in the writing), we could perhaps handle this by writing \emph{For lower triangular matrices, $b$-banded means $b$ total bands, including the main diagonal and $b-1$ bands below the main diagonal. We abuse this notation slightly and call a symmetric banded matrix $X$ $b$-banded when it has a total $2b-1$ digaonals, with $b-1$ bands both above and below the main diagonal.}}

\paragraph{The matrix mechanism and sensitivity}
We utilize the matrix mechanism~\citep{Li2015TheMM}, which, provided a factorization $\bfA = \bfB\bfC$, computes the estimate
\begin{equation}\label{eq:mat_mech}
\widehat{\bfA\obs} = \bfB\left(\bfC\obs + \bfz\right),
\end{equation}
where $\bfz$ is a sample from appropriately scaled isotropic Gaussian noise. The scale is determined by the \emph{sensitivity} of the mapping $\obs \mapsto \bfC \obs$. Following \citet{choquette22multiepoch}, let $\bfN$ be the set of all pairs of neighboring streams $\obs$ and $\deltaset := \left\{ \obs - \obsp \mid (\obs, \obsp) \in \bfN \right\}$ represent the set of all possible deltas between neighboring $\obs, \obsp$. We say a $\deltaset$ \textbf{satisfies the participation schema $\Pi$} if the indices of all nonzero elements in each vector $\contrib \in \deltaset$ correspond to some $\pi \in \Pi$ \btodo{May need to make this more precise for scalar vs vector contributions.}. Critically, for linear queries $\deltaset$ fully captures the sensitivity of the query:
\begin{definition}\label{def:sensopmech}
The \textbf{sensitivity} of the matrix factorization mechanism \cref{eq:mat_mech} is defined as
\begin{equation}\label{eq:sensopnorm}
\sensd(\bfC) = \sup_{(\obs, \obsp) \in \bfN} \| \bfC \obs - \bfC \obsp \|_F
 = \sup_{\contrib \in \deltaset} \| \bfC\contrib \|_F. 
\end{equation}
\end{definition}

\begin{prop}[\citet{choquette22multiepoch}]
The sensitivity of $\bfC$ can be upper bounded by.

$$ \max_{\pi \in \Pi} \sum_{ij} | \bfX_{\pi\pi}^{ij} | $$

where $\bfX = \bfC^{\top} \bfC$.  This upper bound is tight when $\bfX \geq 0$.
\end{prop}

\btodo{We probably cannot avoid introducing $\deltaset^1$ and $\deltaset^\mdim$ and the need for the vector-to-scalar reduction. See e.g. \cref{cor:matrix_to_vector_sens}. Probably best to add that background here.}

\ktodo{One of the interesting observations of thinking in the Grothendieck way is: it's not clear that we must frame the problem as scalar-valued, then invoke vector-to-scalar reduction. E.g. we may recover (and even generalize) many of our results from MEMF from just directly analyzing the trace of interest with vector-valued contributions, including those of the nonnegative subsampled $\bfX$ used here. In general, however, framing the problem as scalar-valued then performing the reduction \emph{must} be significantly sharper in some cases, though these cases still require an oracle or good approximation of an $\ell_\infty$ to $\ell_1$ operator norm. I bring this up here because we might be able to simplify some of our presentation by just immediately going for the $d$-dimensional version, explicitly making our $\bfx$ here matrix-valued and viewing the elements of $\pi \in \Pi$ as rows of a matrix, and bypass the need to discuss the reduction.}

\rnote{+1 to putting some of this stuff in the ICML paper, that way we can avoid this nuance here and just cite the relevant sensitivity theorem from that paper.}

Different factorizations $\bfA = \bfB \bfC$ can have very different performance in practice.  Thus, \citet{denisov22matfact} and \citet{choquette22multiepoch} propose optimizing over the space of factorizations, where the objective function is the expected total squared error on $\bfA$, which can be calculated as the product of sensitivity squared (how much noise do we need to add to answer queries in $\bfC$) and the error profile (how much is this noise magnified to go from answers to $\bfC$ to answers to $\bfA$).

\begin{definition}
The \textbf{error profile} of an encoder $\bfC$ is defined as the magnification in error from estimating answers to the queries in $\bfA$ from noisy answers to the queries in $\bfC$.  It can be calculated as $ \norm{\bfA \bfC^{+}}_F^2$, assuming $\bfA$ is in the rowspace of $\bfC$.
\end{definition}

The expected total squared error is invariant to scaling $\bfC$  by a constant, and hence the optimization problem to minimize the error profile under a sensitivity $1$ constraint is given in \cref{prob:optconvex}.
Note that it is useful to express sensitivity and error profile in terms of $\bfX = \bfC^{\top} \bfC$.  The sensitivity can be expressed as $\max_{\contrib \in \deltaset} \bfu \bfX \bfu^{\top} $, and the error profile can be expressed as $\tr{[\bfA^{\top} \bfA \bfX^{+}]}$.  

\begin{problem} \label{prob:optconvex}
The matrix factorization optimization problem requires solving:
\begin{equation} \label{eq:optconvex}
\begin{aligned}
& \underset{\bfX \in \bfS^{n}_{+}}{\text{minimize}}
& & tr[\bfA^{\top} \bfA \bfX^{-1}] \\
& \text{subject to}
& & \max_{\contrib \in \deltaset} \bfu \bfX \bfu^{\top} \leq 1
\end{aligned}
\end{equation}
And choosing $\bfC$ so that $\bfC^{\top} \bfC = \bfX$, e.g., via Cholesky decomposition.
\end{problem}

This is a convex optimization problem, and as such we can find the global minimum with numerical techniques. The formulation of this problem corresponding to \emph{single participation} has been the most-well studied, where the constraint simplifies to $\diag{(\bfX)} \leq \bf1$ \cite{Li2015TheMM,yuan2016convex,mckenna2021hdmm,denisov22matfact}. It is well-understood how to solve this problem in practice, using both primal and dual algorithms.  The generalization to multiple participations with a fixed epoch order has been studied by \citet{choquette22multiepoch}.  %In this paper, our focus will be on a $\Pi_b$ participation schema, which presents some new technical challenges to overcome.